% Nota històrica: d'on ix l'èpsilon (valencià).

\begin{frame}
\didactatitle{Una breu història de l'èpsilon}

Des de Newton i Leibniz, a finals del segle \textsc{xvii}, es van calcular
límits durant més d'un segle sense definir-los: es raonava amb quantitats
«evanescents» o «infinitament petites», i funcionava, però ningú no sabia dir
amb precisió per què.

\begin{itemize}
  \item \textbf{1817.} Bolzano demostra el teorema dels valors intermedis
        sense recolzar-se en la intuïció geomètrica.
  \item \textbf{1821.} Cauchy, en el seu \emph{Cours d'analyse}, defineix el
        límit i la continuïtat amb paraules: una variable que s'acosta
        indefinidament a un valor fix.
  \item \textbf{Cap a 1860.} Weierstrass, en les seues classes de Berlín, dona
        a la definició la forma amb $\varepsilon$ i $\delta$ que fem servir
        hui.
\end{itemize}

\onlynotes{%
  El que va canviar no va ser el càlcul, sinó la manera d'estar-ne segurs: amb
  $\varepsilon$ i $\delta$, «s'acosta» deixa de ser una imatge i passa a ser
  una desigualtat.
}
\end{frame}
